\subsection{Integral operator limits for a partial character packet}
\label{sec:partial-packet-operators}

\paragraph{Finding A50 (integral limits of a partial-packet projector).}
The unrestricted-quotient supplement to A49 produces one good
signed $p$-orbit of character values. A $p$-adic projector onto
that orbit need not be an integral operator of small coefficient
sum. We give an exact norm bound and a finite module example
which make this distinction precise. Neither assertion constructs
a curve or proves failure of any actual prescribed-point cover.

\paragraph{A quantitative bound for approximating a partial projector.}
Let $n>1$ be prime to a rational prime $p$, put
$U=(\mathbf Z/n)^\times$ and $r=\varphi(n)$, and let
$\varnothing\ne\Omega\subsetneq U$ be a union of orbits of
multiplication by $p$. Fix a primitive $n$th root $\zeta$ in
$\overline{\mathbf Q}_p$. For
\[
 A=\sum_{i=0}^{n-1}a_i\sigma^i\in\mathbf Z[C_n],\qquad
                     d=\sum_i|a_i|,
 \tag{PP.1}
\]
suppose $A(\zeta^u)$ is a $p$-adic unit for at least one
$u\in\Omega$, and, for an integer $a\ge1$,
\[
                    v_p(A(\zeta^u))\ge a
                         \quad(u\in U\setminus\Omega).
 \tag{PP.2}
\]
The normalization is $v_p(p)=1$. Then
\[
                  d\ \ge\ p^{a(r-|\Omega|)/r}.
 \tag{PP.3}
\]
In particular, integral approximations modulo $p^a$ to the
$p$-adic primitive-character projector supported on $\Omega$
cannot have uniformly bounded coefficient sum as $a\to\infty$.

\begin{proof}
The algebraic integer $\alpha=A(\zeta_n)$ in
$\mathbf Q(\zeta_n)$ is nonzero: otherwise every primitive
evaluation would vanish, contrary to the unit assumption.
Its field norm is therefore a nonzero rational integer.
The product over $u\in U$ is that norm. All its factors are
$p$-adically integral, and (PP.2) gives
\[
 v_p\!\left(N_{\mathbf Q(\zeta_n)/\mathbf Q}(\alpha)\right)
                         \ge a(r-|\Omega|).
 \tag{PP.4}
\]
The equality between the product and norm uses all primitive
embeddings at the one fixed $p$-adic embedding; no place is
discarded. Since $p\nmid n$, $p$ is unramified in this cyclotomic
field and the normalization agrees with the usual valuation on
the integer norm. At every complex embedding,
$|A(\zeta_n^u)|\le d$. Hence
\[
 p^{a(r-|\Omega|)}
       \le\left|N_{\mathbf Q(\zeta_n)/\mathbf Q}(\alpha)\right|
       \le d^r,
\]
which proves (PP.3). A projector approximation is one on the
good characters and zero on the others modulo $p^a$, so it
satisfies the stated assumptions. Exact integral isolation of
a nonempty proper part of one rational cyclotomic block is
also impossible: vanishing at one primitive characteristic-zero
root forces divisibility by $\Phi_n$, hence vanishing at all.
\end{proof}

This is a lower bound conditional on the displayed divisibility
requirements, not an estimate of actual geometric torsion
exponents. It leaves open the existence of a useful small
operator for a particular cover and particular collision class.

\paragraph{A complete length-thirty-one spectral calculation.}
Let $G=C_{31}=\langle\sigma\rangle$, let
$K=\mathbf F_{32}=\mathbf F_2[z]/(z^5+z^2+1)$, and put
$\zeta=z$. Its multiplicative order is thirty-one. Define
\[
 \Omega=\{\pm2^j\bmod31:0\le j<5\},\qquad |\Omega|=10.
 \tag{PP.5}
\]
In the binary group algebra $\mathbf F_2[G]$, the entire
subspace whose Fourier transform is supported on $\Omega$
consists of the words
\[
 c(a,b)_i=\operatorname{Tr}_{K/\mathbf F_2}
             (a\zeta^i+b\zeta^{-i}),
       \qquad 0\le i<31,\quad (a,b)\in K^2.
 \tag{PP.6}
\]
It has dimension ten, and its minimum nonzero Hamming weight
is ten. The complete weight distribution is
\[
\begin{array}{c|rrrrrrr}
\text{weight}&0&10&12&14&16&18&20\\
\hline
\text{number of words}&1&31&155&310&217&155&155.
\end{array}
 \tag{PP.7}
\]

\begin{proof}
Expand the trace as a sum of the five $2$-power conjugates.
For the Fourier convention
$\widehat c(k)=\sum_{i=0}^{30}c_i\zeta^{ki}$, character
orthogonality gives
\[
 \widehat{c(a,b)}(-2^j)=a^{2^j},\qquad
 \widehat{c(a,b)}(2^j)=b^{2^j},
 \tag{PP.8}
\]
and zero outside $\Omega$. The two orbits are disjoint:
two has order five modulo thirty-one, and its powers do not
contain minus one. In particular (PP.8) recovers $a,b$ and
proves injectivity of the ten-dimensional binary parameter
space. Conversely a binary word satisfies
$\widehat c(2k)=\widehat c(k)^2$.
On each of the two five-element orbits, one arbitrary element
of $K$ determines all the Fourier values. Thus the entire
binary spectral subspace has dimension at most ten, and
(PP.6) is all of it, not just a subcode.

The finite assertion (PP.7) is verified by enumeration of all
$32^2=1024$ pairs in the retained certificate. It uses exact
arithmetic in the displayed polynomial quotient, checks that
the thirty-one powers of $z$ are all nonzero elements, and
retains every parameter pair, word and weight. The sum of
the displayed multiplicities is $1024$ and every word is distinct.
The pair $a=b=1$ has weight ten. A separate implementation
using five-coefficient polynomial convolution reproduces
the entire distribution and checks the Fourier support of
all ten generators. This is finite exhaustive verification,
not a claimed formal proof-assistant verification.
\end{proof}

\paragraph{A finite module with no sufficiently small annihilator.}
There is a cyclic $\mathbf Z/8[G]$-module $W$, with a generator
$\delta$, such that its good packet $\Omega$ is absent, but
every nonzero $A\in\mathbf Z[G]$ with $A\delta=0$ satisfies
\[
                         \|A\|_1\ge8.
 \tag{PP.9}
\]
Consequently no such nonzero annihilator meets
$31>(\|A\|_1-1)^2$, the small-degree comparison used in
the genus-two integral-operator argument.

\begin{proof}
Let $e_0\in\mathbf F_2[G]$ be the idempotent with Fourier
value one on $\Omega$ and zero elsewhere. It is the word
$c(1,1)$ in (PP.6). Lift its coefficients to integers and
apply the polynomial $f(t)=3t^2-2t^3$ twice in
$R=(\mathbf Z/8)[G]$. Since
\[
 f(t)^2-f(t)=(t^2-t)^2(4t^2-4t-3),
\]
the resulting element $e$ is idempotent modulo eight and
reduces to $e_0$. Define
\[
                        W=(1-e)R,\qquad\delta=1-e.
 \tag{PP.10}
\]
It is cyclic, and multiplication by $e$ kills it. Its reduction
modulo two contains exactly the complement of $\Omega$,
including the trivial character; its rank over $\mathbf Z/8$
is twenty-one. Thus there is a whole missing good signed
packet, not merely one missing character.

Suppose $A\ne0$ and $A\delta=0$.
Write $A=2^bB$ with at least one coefficient of $B$ odd.
If $b\ge3$, then $\|A\|_1\ge8$.
If $b<3$, the equality
$2^bB(1-e)=0$ modulo eight implies, on reduction modulo two,
\[
                         \bar B(1-e_0)=0.
\]
The nonzero binary word $\bar B$ therefore belongs to the
entire spectral subspace of (PP.6), so it has at least ten
nonzero coefficients. Hence
$\|A\|_1=2^b\|B\|_1\ge2^b\cdot10\ge10$.
This proves (PP.9); the bound eight is attained by the scalar
operator $8$, so the unrestricted minimum is exactly eight.
The degree comparison would require $\|A\|_1\le6$, and
therefore fails for every nonzero annihilator in this example.
\end{proof}

For an integral group-ring operator and two disjoint free
deck orbits, the positive part of $A(Q'-P')$ has degree
$\|A\|_1$. Thus (PP.9) addresses the precise coefficient
quantity in that construction. At exponent two, by contrast,
the scalar two already annihilates the entire abstract module;
the binary weight bound alone would only obstruct a lift
which remains nonzero modulo two. Using the exponent-eight
module in (PP.10) avoids that mistaken inference.

This is an abstract finite-module counterexample to an
unqualified small-integral-operator deduction from one good
packet. It is not asserted to be the rational torsion module
of an actual Jacobian, or to arise from a collision of point
sections. It does not refute geometric supply, exclude other
operators on a particular curve, or prove that every successful
argument must use a full faithful conductor. The norm bound
and finite example specify why a $p$-adic packet projector by
itself does not close the remaining step of the A49 supply note.
